\section{Results}
\label{sec:results}

\begin{table*}[ht]
\centering
\begin{tabular}{llccccccc}
\toprule
                                          & \multirow{2}{*}{Method}      & \multirow{2}{*}{\makecell{CMMD }} & \multirow{2}{*}{\makecell{TOPIQ $\uparrow$}} & \multirow{2}{*}{\makecell{VL $\uparrow$}} & \multirow{2}{*}{\makecell{FVD }} & \multirow{2}{*}{\makecell{LipScore $\uparrow$}} 
                                          & \multirow{2}{*}{\makecell{LipLeak }} & \multirow{2}{*}{\makecell{Elo $\uparrow$}} \\ \\ \midrule
\multirow{6}{*}{\STAB{\rotatebox[origin=c]{90}{Reconstruction}}}    & DiffDub~\cite{DiffDub}     &  0.403  &       0.44    & 37.12 &429.07   &   0.34    & -  &   1014  \\         
& IP\_ LAP~\cite{Zhong_2023_CVPR}     & \underline{0.091}   &  \underline{0.49}   &  37.77 &\underline{282.02}    &  0.36    &   -    &  1007     \\
& Diff2Lip~\cite{diff2lip}   &   0.225     &   0.48    & 35.84 & 555.08   &   0.49    &  -   &  886    \\
                                          & TalkLip~\cite{talklip}   & 0.230   &  0.39     & 29.07  &608.92   & \textbf{0.58}   &    -   & 920       \\ 
                                          & LatentSync~\cite{latentsync}   & 0.319   &  0.41     & \underline{45.23}  &343.90   & \underline{0.52}   &    -   &    \underline{1052}    \\ 
                                          & \cellcolor{cvprblue!40}KeySync &  \cellcolor{cvprblue!40}\textbf{0.064} &  \cellcolor{cvprblue!40}\textbf{0.58}  &  \cellcolor{cvprblue!40}\textbf{70.32} & \cellcolor{cvprblue!40}\textbf{191.21}  &  \cellcolor{cvprblue!40}0.46    & \cellcolor{cvprblue!40}- & \cellcolor{cvprblue!40}\textbf{1120} \\ \hline
\multicolumn{1}{c}{\multirow{6}{*}{\STAB{\rotatebox[origin=c]{90}{Cross-sync}}}} 

& DiffDub~\cite{DiffDub}         &  0.408  &    0.44   & 37.05 &420.66     & \underline{0.34}    & 0.56     &  947       \\
\multicolumn{1}{c}{}                      & IP\_\ LAP~\cite{Zhong_2023_CVPR}    &  \underline{0.093}  &   0.49  & 35.32 &\underline{294.66}   &  0.17    &  0.28     &     1031      \\
\multicolumn{1}{c}{}                      & Diff2Lip~\cite{diff2lip}   & 0.231  & \underline{0.48}   &33.97 &601.68  &  0.16    & \underline{0.25}    &       878        \\

\multicolumn{1}{c}{}                      & TalkLip~\cite{talklip}    & 0.201  &  0.42  & 24.80  &704.93   &  0.30  &  0.66      &    911    \\
& LatentSync~\cite{latentsync}   & 0.325  &  0.41     & \underline{45.95}  &361.57   & 0.14   &    0.33   &   \underline{1086}    \\ 
\multicolumn{1}{c}{}                      & \cellcolor{cvprblue!40}KeySync &  \cellcolor{cvprblue!40}\textbf{0.070}
 &  \cellcolor{cvprblue!40}\textbf{0.58}  &  \cellcolor{cvprblue!40}\textbf{73.04} & \cellcolor{cvprblue!40}\textbf{206.32}  &  \cellcolor{cvprblue!40}\textbf{0.48}    & \cellcolor{cvprblue!40}\textbf{0.16}   & \cellcolor{cvprblue!40}\textbf{1145}  \\ 
                                         \bottomrule
\end{tabular}

\caption{\textbf{Quantitative comparison with other works} on reconstruction and cross-synchronization performance. The best results are highlighted in \textbf{bold}, while the second-best results are \underline{underlined}. All metrics are described in Section~\ref{sec:metrics}.
}
\label{tab:metrics_comparison}
\end{table*}

In this section, we present a comprehensive evaluation of our model’s performance against baselines, along with several ablations to assess the impact of key components.

\subsection{Comparison with other works} \label{sec:quantitative}
\paragraph{Quantitative analysis.}

We evaluate our method alongside five competing approaches in Table~\ref{tab:metrics_comparison}. The evaluation is conducted in two settings: reconstruction, where videos are generated using the same audio as in the original video, and cross-sync, where the audio is taken from a different video. The latter is particularly relevant as it better reflects real-world applications such as automated dubbing, where the driving audio is typically not aligned with the input video. 

KeySync consistently outperforms other methods in visual quality and temporal consistency for both reconstruction and cross-sync tasks, as reflected in the higher VL scores and lower FVD. Regarding lip synchronization, while most methods experience a drop in LipScore during cross-sync, our approach maintains consistent performance, as evidenced by similar LipScores in both reconstruction and cross-sync. In the reconstruction setting, some methods achieve a higher LipScore, but this is primarily due to expression leakage. This hypothesis is also supported by the high LipLeak scores exhibited by these models. An interesting observation arises with DiffDub: in the cross-sync setting, the model generates random lip movements, which LipScore mistakenly interprets as improved synchronization, but LipLeak exposes as leakage. Our method is also preferred by human evaluators in both settings, as shown by the higher Elo rankings, highlighting the effectiveness of our approach according to human perception.

\paragraph{Qualitative analysis.}

Figure~\ref{fig:qualitative} presents results from all models for a cross-sync example. We see that KeySync more accurately follows the lip movements corresponding to the input audio. While LatentSync and Diff2Lip also appear to align somewhat with the target lip movements, they fail to generate certain vocalizations correctly and exhibit visual artifacts (highlighted on the figure via red squares and arrows, respectively), limiting their practical usability. Additionally, we notice that most methods produce minimal lip movement or fail to open the mouth sufficiently. This can be attributed to expression leakage, where the model receives conflicting signals from the original video and the new audio, making it difficult to generate a coherent and natural-looking mouth region.

\paragraph{Leakage.}

As discussed in Section~\ref{sec:metrics}, we compute LipLeak by generating a video using a silent audio input. Since the audio contains no speech, the mouth should remain closed throughout the entire duration. However, in practice, we observe that this is not always the case, as the non-silent expressions in the input video can leak into the generated video. Figure~\ref{fig:leakage} presents examples of such generated videos alongside the original input video and silent audio. We observe that all methods, except ours and Diff2Lip, exhibit several frames where the mouth is open (highlighted by red squares around the mouths) due to expression leakage from the input video. While Diff2Lip manages to keep the mouth closed, we notice significant blending artifacts across all frames, highlighting the model’s struggle to generate silent frames when the original video contains speech. Additionally, in Figure~\ref{fig:leakage_plot}, we visualize the mouth aspect ratio (MAR) of all methods over time. Since MAR is a ratio, it remains independent of scale, ensuring a fair comparison across different methods. The results show that baseline methods frequently exceed the MAR threshold for an open mouth, indicating persistent leakage. In contrast, KeySync consistently maintains a MAR below the threshold, demonstrating its effectiveness in preventing unwanted lip movements when no speech is present.


\paragraph{Occlusion handling.}
\begin{figure}
    \centering
    \includegraphics[width=1.0\linewidth]{Images/quali_comp_v1.pdf}
    \caption{\textbf{Qualitative comparison with other works.} The top row (“Target lips”) shows lip movements corresponding directly to the provided audio input, and can therefore be seen as the target for the lips in the generated videos.}
    \label{fig:qualitative}
\end{figure}
\begin{figure}
    \centering
    \includegraphics[width=1.0\linewidth]{Images/leakage_v2.pdf}
    \caption{\textbf{Qualitative leakage comparison.} We condition the models on silent audio and non-silent video (first row).}
    \label{fig:leakage}
\end{figure}
\begin{figure}
    \centering
    \includegraphics[width=\linewidth]{Images/mar_comparison.pdf}
    \caption{\textbf{MAR over time.} If MAR exceeds the threshold, the mouth is considered open, indicating leakage.}
    \label{fig:leakage_plot}
\end{figure}
\begin{figure*}
    \centering
    \includegraphics[width=\textwidth]{Images/occlusion_quali_v1.pdf}
    \caption{\textbf{Occlusion handling comparison.} We present qualitative results on the left and quantitative results on the right.}
    \label{fig:occlusions}
\end{figure*}

We propose a technique to handle occlusions at inference time, as defined in Section~\ref{sec:occlusions} and illustrated in Figure~\ref{fig:architecture}. The effectiveness of our approach is demonstrated in Figure~\ref{fig:occlusions}. On the left side, we observe that without proper occlusion handling, the model fails to reconstruct the occluded regions accurately, leading to unwanted artifacts, particularly around the hand. In contrast, with our proposed method, the hand is correctly represented while maintaining lip synchronization. This improvement is further highlighted in the right part of the figure, where we visualize the mean absolute error between the generated video and the ground-truth. Without occlusion handling, we observe noticeable error spikes during occlusions. 
% On the other hand, our method, which leverages a large context field, not only mitigates these errors but also enhances the reconstruction both before and after the occlusion. 
Additional details are provided in the supplementary material.

\subsection{Ablation studies}

%To effectively evaluate the key aspects of KeySync, we conduct a series of ablation studies.

\paragraph{Architecture.}

We evaluate two alternate versions of our pipeline and present the results in Table~\ref{tab:architecture_ablation}. The first approach replaces our two-stage pipeline with a single-stage model, where the network directly generates a sequence of frames, and a longer video is produced by concatenating multiple sequences with a one-frame overlap to ensure temporal consistency. The second approach keeps the two-stage logic,  but generates keyframes using an image-based model rather than producing them all simultaneously with temporal layers. We find that the one-stage model achieves reasonable visual quality according to CMMD, but sees a sharp decline in FVD and LipScore, highlighting the importance of our keyframe interpolation technique for generating smooth, well-synchronized lip movements. Similarly, without the temporal layers, the interpolation model struggles to maintain coherence, leading to significant degradation across all three metrics.

\paragraph{Audio encoder.}
\begin{table}[t]
\centering
\resizebox{\columnwidth}{!}{
\begin{tabular}{ccccc}
\toprule
Two-stage & Temp. layers & CMMD & FVD  & LipScore $\uparrow$ \\ \midrule
  \xmark & \cmark      &    \underline{0.085}              & \underline{395.45} & 0.32   \\
  \cmark     & \xmark  & 0.142             & 618.27 & \underline{0.39} \\
   \rowcolor{Gray!40} \cmark     & \cmark  & \textbf{0.070}                  &  \textbf{206.32} &  \textbf{0.48} \\
\bottomrule
\end{tabular}}
\caption{\textbf{Architecture ablation} in the cross-sync setting.}
\label{tab:architecture_ablation}
\end{table}
\begin{table}[t]
\centering
% \resizebox{\columnwidth}{!}{%
\begin{tabular}{lcccc}
\toprule
Audio backbone     & FVD  & LipScore $\uparrow$ & LipLeak  \\ \midrule
Whisper~\cite{whisper}    &  207.41  &  0.47   &  0.18  \\
Wav2vec2~\cite{wav2vec2}         & \textbf{201.13}    & 0.45  & 0.19  \\ 
 WavLM~\cite{wavlm}    &  218.08   & \textbf{0.48}  & 0.17  \\
\rowcolor{Gray!40} HuBERT~\cite{hubert}        & \underline{206.32}    & \textbf{0.48}  & \textbf{0.16}  \\
\bottomrule
\end{tabular}
\caption{\textbf{Audio encoder ablation} in the cross-sync setting.}
\label{tab:abl_audio}
\end{table}
\begin{table}[t]
\centering
\resizebox{\columnwidth}{!}{%
\begin{tabular}{lcccc}
\toprule
Mask     &        CMMD & FVD & LipScore $\uparrow$ & LipLeak  \\ \midrule
Mouth-only      &  0.077  &  \underline{200.71}   & 0.23  &  0.38 \\
Full lower-face   &  0.743   &  219.96   & 0.35  & \underline{0.28}  \\
Ours (nose-level)   &  \underline{0.071}   &  \textbf{199.39}   &  \underline{0.34} & 0.35 \\
\rowcolor{Gray!40} Ours &  \textbf{0.070}   &  206.32   & \textbf{0.48}  & \textbf{0.16}  \\
\bottomrule
\end{tabular}}
\caption{\textbf{Mask ablation} in the cross-sync setting.}
\label{tab:abl_mask}
\end{table}

We also investigate the impact of different audio encoders on the generated videos, as shown in Table~\ref{tab:abl_audio}. We see that Wav2vec2~\cite{wav2vec2} produces marginally higher video quality, as indicated by its lower FVD score. However, this comes at the expense of lip synchronization, as reflected in its lower LipScore. With WavLM~\cite{wavlm}, we achieve a LipScore comparable to HuBERT~\cite{hubert}, but at the cost of worse video quality. In contrast, HuBERT not only maintains a strong LipScore but also achieves the lowest LipLeak, indicating its effectiveness in mitigating expression leakage. Based on these findings, we select HuBERT as our default audio encoder.

\paragraph{Mask.}
\begin{figure}
    \centering
    \includegraphics[width=\linewidth]{Images/masking_tech_v1.pdf}
    \caption{\textbf{Examples of different masking techniques.}}
    \label{fig:masking}
\end{figure}

Finally, we investigate the impact of different masking techniques (illustrated in Figure~\ref{fig:masking}) in Table~\ref{tab:abl_mask}. Using a mouth-only mask yields better video quality since it minimizes facial obstruction. However, this approach suffers from severe leakage, as indicated by its low LipScore and high LipLeak. This occurs because the model can infer mouth movements from the mask’s position over time rather than learning proper synchronization with the driving audio. Conversely, masking the entire lower face effectively reduces leakage, but severely harms image and video quality, as the model is forced to reconstruct unrelated background elements. Our proposed box-style masking offers a balanced trade-off between these two extremes, achieving the best overall performance. Additionally, we demonstrate that extending the mask up to the eye region is crucial in maximizing LipScore, as the cheeks convey important cues about mouth movements and can therefore cause leakage. 